% Cover letter -- General Relativity and Gravitation (source of the PDF; content
% mirrors COVER_PHOTON_ARC_GRG.md). Compile: pdflatex COVER_PHOTON_ARC_GRG.tex
\documentclass[11pt]{article}
\usepackage[a4paper,margin=1.1in]{geometry}
\usepackage[T1]{fontenc}
\usepackage[utf8]{inputenc}
\usepackage{amsmath,amssymb}
\usepackage[hidelinks]{hyperref}
\usepackage{parskip}
\usepackage{enumitem}
\pagestyle{empty}

\begin{document}

\begin{flushright}
Miqueias Alves Mendes\\
Independent Researcher\\
Ibiapina, Cear\'a, Brazil\\
\href{mailto:mikalvesm@gmail.com}{mikalvesm@gmail.com}\\[4pt]
\today
\end{flushright}

\vspace{6pt}
\noindent To the Editors of \emph{General Relativity and Gravitation},

\vspace{6pt}
I submit for your consideration the manuscript \textbf{``An indefinite-signature
Benincasa--Dowker gauge operator for causal sets: electric causal diamonds, an
$H^2$ curvature law, and the obstruction map for an emergent photon.''}

\textbf{What the paper delivers.} Causal set theory carries a smeared,
Lorentz-invariant d'Alembertian for \emph{scalar} fields (the Benincasa--Dowker
line), but no analogous operator for a gauge field. The manuscript contributes
three portable objects:

\begin{enumerate}[leftmargin=1.6em]
\item \textbf{A new operator.} An indefinite-signature ($E^2-B^2$)
Benincasa--Dowker gauge operator---to my knowledge the first BD-type operator for
a gauge $1$-form---in which each causal-diamond plaquette enters with a sign set
by its area bivector (timelike/electric vs.\ spacelike/magnetic). It is
gauge-invariant to $5.7\times10^{-16}$ and reproduces free Maxwell on a $4$D
lattice ($\omega=ck$ with fitted $c=1.05$; $c$ is never input). Applied to the
Poisson causal set it returns a sharp structural verdict: causal diamonds are
$100\%$ electric (magnetic-cell fraction $0.0000$ exactly, Wilson bound
$6\times10^{-5}$, stable for $N=200$--$2000$)---the magnetic sector a Lorentzian
signature requires is structurally empty.

\item \textbf{A proposition.} The non-locality that obstructs every bare
connection on the sprinkling is proved structural: by a Campbell--Mecke (Palm)
argument, \emph{every} Poincar\'e-invariant connection rule---pairwise or
depending on any invariantly defined neighbourhood configuration---has divergent
expected valency, because the boost orbit at fixed proper time is the non-compact
hyperboloid. The familiar non-locality of causal-set links is the special case of
the covering relation; the proposition shows it holds for every invariant rule,
and that the natural candidate exception (the Alexandrov-interval abundance
$N(i,j)$) is the same variable in disguise.

\item \textbf{A map.} A pre-registered chain of measurements that relocates the
obstruction to an emergent photon step by step: the orientation Goldstone modes
are internal scalars, not a gauge vector; the link-$U(1)$ sector is not
permanently confining but its Coulomb phase is not certifiable on the non-local
order; diamond height does not supply the missing spacelike cell, and an ordering
coupling does not amplify it---but spatial curvature does, with a clean law
$\Delta f_B\propto(1/\hat R)^{1.73}\approx H^2$ ($R^2=0.997$). An $O(10\%)$
magnetic sector then requires sub-Planckian curvature: the mechanism is real but
Planck-scale, and does not explain the photons of the observable low-curvature
universe. Two further low-curvature levers (future-cone cells; a genuine
Bianchi-I anisotropy on the causal order) are exhausted as well, closing the
low-curvature map. The remaining obstruction is stated exactly: a
Lorentz-invariant spacelike $2$-cell on the causal order at low curvature.
\end{enumerate}

\textbf{What is claimed, carefully.} No photon is claimed, and no impossibility
either. The indefinite operator has no continuum theorem of its own (it is
validated numerically against the known lattice answer), so its causet verdict is
read as a structural statement about the causal-diamond $2$-complex, not a proven
continuum limit; the identification of the discreteness scale with the Planck
scale is declared as an external assumption. Every number is produced by code
that passes an anti-circularity guard forbidding any relativistic input in the
data-generating modules.

\textbf{Why \emph{General Relativity and Gravitation}.} The work is causal set
theory and discrete gauge structure on a quantum-gravity substrate: a new
operator extending the Benincasa--Dowker line from scalars to gauge $1$-forms, a
proposition about Lorentz-invariant rules on sprinklings, and the curvature
dependence of an emergent magnetic sector. It tells the community precisely where
an emergent photon in causal set theory must be sought and why the natural routes
fail, with a reusable operator and a theorem as byproducts---squarely within the
journal's scope in classical and quantum gravitation.

\emph{Note on manuscript class.} The source is currently prepared in REVTeX; upon
acceptance in principle I will supply a version in the Springer class to meet the
journal's formatting requirements.

The manuscript is original, is not under consideration elsewhere, and reports the
author's own simulations; the figures and tables are generated from the same runs
that produce the quoted numbers. I declare no conflicts of interest.

Thank you for considering the manuscript.

\vspace{10pt}
Sincerely,

Miqueias Alves Mendes
\end{document}
